% T6 — «Минимал»: без рамок, тонкая линия слева, sans-serif, серый акцент.
% Профиль: математика, простые числовые задачи. 7 задач, 1 колонка.
\documentclass[a4paper,11pt]{article}

\usepackage{../../../../Lessons/_templates/preamble_worksheet}
\usepackage{../_shared/worksheet_check}

% --- Палитра: благородный серо-стальной акцент ---
\definecolor{wcprimary}{HTML}{4A5568}
\definecolor{wcprimaryLight}{HTML}{CBD5E0}
\definecolor{wcprimaryBG}{HTML}{F7FAFC}

% Глобально включаем sans-serif на весь документ
\renewcommand{\familydefault}{\sfdefault}

% --- Заголовок: минимал, тонкая линия снизу ---
\renewcommand{\WorksheetTitle}[1]{%
  \par\noindent
  {\fontsize{20pt}{24pt}\selectfont\bfseries\color{wcprimary}#1}\par
  \vspace{1mm}{\color{wcprimaryLight}\rule{\linewidth}{0.4pt}}\par\vspace{4mm}}

% --- TaskBox: только тонкая вертикальная линия слева + номер ---
\renewcommand{\TaskBox}[2]{%
  \par\noindent\begin{tikzpicture}
    \node[anchor=north west, inner sep=0pt, text width=\linewidth] (B) at (0,0) {%
      \begin{minipage}[t]{\linewidth}
        \noindent\hspace*{4mm}{\small\bfseries\color{wcprimary}\##1}\par
        \vspace{1mm}\hspace*{4mm}\begin{minipage}{\dimexpr\linewidth-4mm\relax}#2\end{minipage}
      \end{minipage}};
    \draw[color=wcprimary, line width=1.4pt]
      ([yshift=-0.5mm]B.north west) -- ([yshift=0.5mm]B.south west);
  \end{tikzpicture}\par\vspace{4mm}}

\SetAnswerFieldSize{28mm}{8mm}

\begin{document}

\WorksheetTitle{Производная и её применение}
{\small\itshape\color{wcprimary!80}ЕГЭ профиль, задание №7 \quad·\quad Класс 11 \quad·\quad T6 \quad·\quad ID: T6-derivative-2026-05-26}\par\vspace{3mm}
\FirstPageHeader

\TaskBox{1}{%
  Найдите значение производной функции $f(x) = 3x^2 - 5x + 4$ в точке $x_0 = 2$.\par
  \vspace{1mm}\hfill\textit{\small Ответ:}~\AnswerField{1}%
}

\TaskBox{2}{%
  Функция $y = f(x)$ определена на интервале $(-6;\,8)$. На рисунке мысленно представьте график её производной с тремя нулями: $-3$, $1$ и $5$. Найдите количество точек экстремума функции $f(x)$.\par
  \vspace{1mm}\hfill\textit{\small Ответ:}~\AnswerField{2}%
}

\TaskBox{3}{%
  Прямая $y = 5x - 7$ параллельна касательной к графику функции $y = x^2 + 6x - 8$. Найдите абсциссу точки касания.\par
  \vspace{1mm}\hfill\textit{\small Ответ:}~\AnswerField{3}%
}

\TaskBox{4}{%
  Материальная точка движется прямолинейно по закону $x(t) = \tfrac{1}{3}t^3 - 2t^2 + 8t - 3$. В какой момент времени (в секундах) её скорость была равна $3$~м/с?\par
  \vspace{1mm}\hfill\textit{\small Ответ:}~\AnswerField{4}%
}

\TaskBox{5}{%
  Найдите точку минимума функции $y = x^3 - 3x^2 + 5$.\par
  \vspace{1mm}\hfill\textit{\small Ответ:}~\AnswerField{5}%
}

\TaskBox{6}{%
  Найдите наибольшее значение функции $y = -x^2 + 4x + 7$ на отрезке $[0;\,3]$.\par
  \vspace{1mm}\hfill\textit{\small Ответ:}~\AnswerField{6}%
}

\TaskBox{7}{%
  Найдите наименьшее значение функции $y = x^3 - 6x^2 + 9x + 1$ на отрезке $[0;\,5]$.\par
  \vspace{1mm}\hfill\textit{\small Ответ:}~\AnswerField{7}%
}

\WorksheetQR{T6-derivative/L1/v1}

\end{document}
